% Sistemas 2x2 bien y mal condicionados: el efecto de perturbar el segundo término independiente.
\documentclass[tikz,border=4pt]{standalone}
\usepackage{resonancia}
\begin{document}
\begin{tikzpicture}[x=2.3cm, y=2.3cm]
  % --- (a) bien condicionado: A = [1 1; 1 -1], cond = 1
  \begin{scope}
    \draw[eje] (-0.1,0) -- (2.1,0) node[right] {$x$};
    \draw[eje] (0,-0.1) -- (0,2.1) node[above] {$y$};
    \draw[medida, line width=1.2pt] (0,2) -- (2,0);
    \draw[color=rNaranja, line width=1.2pt] (0,0) -- (2,2);
    \draw[color=rNaranja, dashed, line width=0.8pt] (0.1,0) -- (2,1.9);
    \fill[rTinta] (1,1) circle (0.035);
    \fill[rNaranja] (1.05,0.95) circle (0.035);
    \node[etiqueta, anchor=east] at (0.9,1.0) {solución $(1;\,1)$};
    \node[etiqueta, color=rNaranja, anchor=west] at (1.25,1.0) {perturbada $(1{,}05;\,0{,}95)$};
    \node[etiqueta] at (1,-0.28) {(a) $\kappa_2(\mathbf A)=1$};
  \end{scope}
  % --- (b) mal condicionado: A = [1 1; 1 1,1], cond ≈ 42
  \begin{scope}[shift={(3.2,0)}]
    \draw[eje] (-0.1,0) -- (2.1,0) node[right] {$x$};
    \draw[eje] (0,-0.1) -- (0,2.1) node[above] {$y$};
    \draw[medida, line width=1.2pt] (0,2) -- (2,0);
    \draw[color=rNaranja, line width=1.2pt] (0,1.909) -- (2,0.091);
    \draw[color=rNaranja, dashed, line width=0.8pt] (0,1.955) -- (2,0.136);
    \fill[rTinta] (1,1) circle (0.035);
    \fill[rNaranja] (0.5,1.5) circle (0.035);
    \draw[punteada] (1,1) -- (1.45,1.45);
    \draw[punteada, color=rNaranja] (0.5,1.5) -- (0.8,1.85);
    \node[etiqueta, anchor=south west] at (1.4,1.4) {solución $(1;\,1)$};
    \node[etiqueta, color=rNaranja, anchor=south west] at (0.75,1.8) {perturbada $(0{,}5;\,1{,}5)$};
    \node[etiqueta] at (1,-0.28) {(b) $\kappa_2(\mathbf A)\approx 42$};
  \end{scope}
\end{tikzpicture}
\end{document}
